% ============================================================
\section{执行最小路径电阻分析}
\subsection*{Performing Minimum Path Resistance Analysis}

\begin{noteBox}
最小路径电阻分析在 RedHawk Fusion 与 RedHawk-SC Fusion 分析流程中均受支持。
\end{noteBox}

执行最小路径电阻分析可快速检测设计中电源与地网络的薄弱点。节点的最小路径电阻值是到理想电压源位置（tap）的最小阻性路径上的电阻值。理想电压源可以是电源或地 pin、用户定义的 tap 或 package。由于该值仅表示单条路径上的电阻，它不一定等于从该节点到电源或地的有效电阻；相反，它表示节点到电源或地有效电阻的上界。

要执行最小路径电阻分析，运行 \cmd{analyze\_rail} \opt{-min\_path\_resistance}。必须用 \opt{-nets} 指定要分析的电源与地 net。

\begin{noteBox}
在 RedHawk-SC 轨分析流程中，必须同时指定 \opt{-voltage\_drop} 与 \opt{-min\_path\_resistance} 以计算最小路径电阻。
\end{noteBox}

可在压降分析之前或期间执行最小路径电阻分析。分析完成后，工具将最小路径电阻值与压降结果写入同一轨结果。

同时执行压降分析与最小路径电阻计算：
\begin{lstlisting}
fc_shell> analyze_rail -nets {VDD VSS} \
   -min_path_resistance -voltage_drop static
\end{lstlisting}

仅执行最小路径电阻分析：
\begin{lstlisting}
fc_shell> analyze_rail -nets {VDD VSS} -min_path_resistance
\end{lstlisting}

\subsection{查看最小路径电阻分析结果}
\subsubsection*{Viewing Minimum Path Resistance Analysis Results}

分析完成后，运行 \cmd{report\_rail\_result} \opt{-type} \texttt{minimum\_path\_resistance}，以显示最小路径电阻地图、将结果写出到输出文件，或查询单元实例或 pin 上的属性。

下例将 VDD net 的计算最小路径电阻写出到 \texttt{minres.rpt}：
\begin{lstlisting}
fc_shell> open_rail_result
fc_shell> report_rail_result -type minimum_path_resistance \
   -supply_nets { vdd } minres.rpt
fc_shell> sh cat minres.rpt
FEED0_14/vbp 518.668
INV61/vdd 158.268
ARCR0_1/vbp 145.582
FEED0_5/vdd 156.228
INV11/vbp 518.669
\end{lstlisting}

要为单元实例写出详细最小电阻路径报告，运行 \cmd{report\_rail\_minimum\_path} 并指定要报告的 net 与单元。要在 GUI 中查看该最小电阻路径的逐段电阻，使用 \cmd{report\_rail\_minimum\_path} 的 \opt{-error\_cell}，将最小路径电阻路径上的几何保存到错误单元。随后可在错误浏览器中打开该错误单元，检查从单元 pin 到最近 tap 的最小电阻路径并追踪途经的每个几何。

下例追踪最小电阻路径上的全部几何并写入名为 \texttt{test.err} 的错误单元：
\begin{lstlisting}
fc_shell> open_rail_result
fc_shell> report_rail_minimum_path -cell U19 -net VDD -error_cell
\end{lstlisting}

如图~\ref{fig:mpr-err} 所示，错误浏览器中按层对每个几何分类。单击几何可查看 IR drop、电阻、包围盒等详细信息。在 GUI 中，追踪路径在设计布局中以黄色高亮（见图~\ref{fig:mpr-trace}）。

\begin{figure}[htbp]
\centering
\fbox{\parbox{0.85\textwidth}{\centering\vspace{1.5cm}
\textit{（原书 Figure 206：Displaying Error Cells in the Error Browser）}\\[0.3em]
\small 请参阅原版 PDF 插图。
\vspace{1.5cm}}}
\caption{在错误浏览器中显示错误单元}
\label{fig:mpr-err}
\end{figure}

\begin{figure}[htbp]
\centering
\fbox{\parbox{0.85\textwidth}{\centering\vspace{1.5cm}
\textit{（原书 Figure 207：Traced Segment is Highlighted in the Design Layout）}\\[0.3em]
\small 请参阅原版 PDF 插图。
\vspace{1.5cm}}}
\caption{追踪线段在设计布局中高亮}
\label{fig:mpr-trace}
\end{figure}

\begin{seeAlsoBox}
在 GUI 中显示地图；写出分析与检查报告；查询属性
\end{seeAlsoBox}

\subsection{使用鼠标工具追踪最小路径电阻}
\subsubsection*{Tracing Minimum Path Resistance Using the Mouse Tool}

使用最小路径电阻（MPR）鼠标工具，可交互式追踪从指定电源或地 net 到设计中边界节点的最小电阻路径。

要用 MPR 交互式追踪工具追踪计算得到的最小路径电阻：
\begin{enumerate}
  \item 在 GUI 中选择 Task $>$ Rail，然后从任务列表选择 Analysis Tools，打开 Task Assistant 窗口（见图~\ref{fig:mpr-tool}）。
  \item 在 Task Assistant - Rail 窗口中单击 MPR Mouse Tool。
  \item 指定要分析的 net，单击 OK。
\end{enumerate}

图~\ref{fig:mpr-trace-map} 所示的 MPR 交互式鼠标工具提供下列功能：
\begin{itemize}
  \item 放大/缩小以查看追踪
  \item 两种操作模式：Add 模式在地图中追踪多条路径；Replace 模式每次在地图中追踪一条路径
  \item 使用标准查询工具显示追踪信息
  \item 选择区域以一次追踪多个单元
  \item 用脚本记录交互操作
\end{itemize}

\begin{figure}[htbp]
\centering
\fbox{\parbox{0.85\textwidth}{\centering\vspace{1.5cm}
\textit{（原书 Figure 208：Invoking the MPR Mouse Tool）}\\[0.3em]
\small 请参阅原版 PDF 插图。
\vspace{1.5cm}}}
\caption{调用 MPR 鼠标工具}
\label{fig:mpr-tool}
\end{figure}

\begin{figure}[htbp]
\centering
\fbox{\parbox{0.85\textwidth}{\centering\vspace{1.5cm}
\textit{（原书 Figure 209：Tracing the Path With the Least Resistance）}\\[0.3em]
\small 可单击隐藏地图中全部对象，然后选择 Cell。请参阅原版 PDF 插图。
\vspace{1.5cm}}}
\caption{追踪最小电阻路径}
\label{fig:mpr-trace-map}
\end{figure}

% ============================================================
\section{执行有效电阻分析}
\subsection*{Performing Effective Resistance Analysis}

\begin{noteBox}
有效电阻分析在 RedHawk Fusion 与 RedHawk-SC Fusion 分析流程中均受支持。
\end{noteBox}

要计算指定电源与地 net 的有效电阻，使用 \cmd{analyze\_rail} \opt{-effective\_resistance}。

要指定用于有效电阻计算的单元实例列表，在文本文件中描述目标实例，然后用 \appopt{rail.effective\_resistance\_instance\_file} 指定该文本文件。

例如：
\begin{lstlisting}
fc_shell> sh cat cell_list.txt
u_GS_PRO_65
u_GS_PRO_64
ChipTop_U_compressor_mode/U3

fc_shell> set_app_options \
   -name rail.effective_resistance_instance_file \
   -value cell_list.txt
fc_shell> analyze_rail -nets {VDD VSS} -effective_resistance
\end{lstlisting}

可在压降分析之前或期间执行有效电阻分析。分析完成后，工具将有效电阻值与压降结果写入同一轨结果。

\begin{noteBox}
在 RedHawk-SC 轨分析流程中，必须同时指定 \opt{-effective\_resistance} 与 \opt{-voltage\_drop} 以计算有效电阻。
\end{noteBox}

计算完成后，工具将结果保存在工作目录中。要将有效电阻值写出到文本文件，使用 \cmd{report\_rail\_result} \opt{-type} \texttt{effective\_resistance}。要查询电阻值，使用 \cmd{get\_attribute}。

例如：
\begin{lstlisting}
fc_shell> get_attribute [get_pins u_GS_PRO_65/VDD] \
   effective_resistance
\end{lstlisting}

\begin{seeAlsoBox}
指定 RedHawk 与 RedHawk-SC 工作目录；写出分析与检查报告
\end{seeAlsoBox}

% ============================================================
\section{执行分布式 RedHawk Fusion 轨分析}
\subsection*{Performing Distributed RedHawk Fusion Rail Analysis}

\begin{noteBox}
分布式轨分析功能仅在 RedHawk Fusion 分析流程中受支持。
\end{noteBox}

默认情况下，\cmd{analyze\_rail} 使用单进程执行 RedHawk 轨分析。为降低内存使用与总运行时间，可使用分布式处理方法，在不同于运行 Fusion Compiler 工具的机器上运行 RedHawk Fusion。

可在同一 Fusion Compiler 会话中向目标 farm 机器提交多个 RedHawk 作业并并行运行。从而仅用一个 Fusion Compiler 许可证即可运行多次 RedHawk 分析。每个 RedHawk 作业需要一个 \texttt{SNPS\_INDESIGN\_RH\_RAIL} 许可证密钥。

要检查生成的分析结果，先运行 \cmd{open\_rail\_result} 加载保存在指定目录中的结果。随后可在地图中调查问题区域，并在文本文件或错误视图中检查违例。

一次只能打开一个结果上下文。要加载另一次分析运行生成的轨上下文，先用 \cmd{close\_rail\_result} 关闭当前轨结果，再打开期望的轨结果（见「在 GUI 中显示地图」）。

\subsection{许可证要求}
\subsubsection*{License Requirement}

默认情况下，所需许可证数量等于分布式处理模式中用 \opt{-multiple\_script\_files} 指定的脚本文件或配置数量。若运行 \cmd{analyze\_rail} 时可用许可证数少于脚本文件数，命令以错误消息停止。

要在 \cmd{analyze\_rail} 运行期间持续检查可用的 \texttt{SNPS\_INDESIGN\_RH\_RAIL} 许可证，将 \appopt{rail.license\_checkout\_mode} 设为 \texttt{wait}。若工具检测到许可证变为可用，会立即重用或检出许可证并向 farm 机器提交脚本。默认 \appopt{rail.license\_checkout\_mode} 为 \texttt{nowait}。

分析结束后，运行 \cmd{remove\_licenses} 释放 \texttt{SNPS\_INDESIGN\_RH\_RAIL} 许可证：
\begin{lstlisting}
fc_shell> remove_licenses SNPS_INDESIGN_RH_RAIL
\end{lstlisting}

如图~\ref{fig:lic-wait} 所示，\appopt{rail.license\_checkout\_mode} 设为 \texttt{wait}。有六个脚本要提交到 farm 机器，但仅有两个可用 \texttt{SNPS\_INDESIGN\_RH\_RAIL} 许可证。此时工具先检出两个许可证并等待其它许可证可用。当某脚本完成作业后，工具立即重用其许可证以提交其它脚本。

\begin{figure}[htbp]
\centering
\fbox{\parbox{0.85\textwidth}{\centering\vspace{1.5cm}
\textit{（原书 Figure 210：Waiting for Available Licenses During Distributed RedHawk Fusion Run）}\\[0.3em]
\small 请参阅原版 PDF 插图。
\vspace{1.5cm}}}
\caption{分布式 RedHawk Fusion 运行期间等待可用许可证}
\label{fig:lic-wait}
\end{figure}

\subsection{启用分布式处理的步骤}
\subsubsection*{Steps to Enable Distributed Processing}

要为 RedHawk Fusion 启用分布式处理：

\begin{enumerate}
  \item 使用 \cmd{set\_host\_options} 定义分布式处理配置。可指定下列一个或多个选项：

\begin{tabularx}{\textwidth}{@{}l X@{}}
\toprule
\textbf{选项} & \textbf{说明} \\
\midrule
\opt{-submit\_protocol} & 指定提交作业的协议，例如 LSF、GRD 或 SGE。注意：提交作业前，运行 \texttt{rsh} 检查目标 farm 机器是否支持从远程机器接受分布式处理的 rsh 协议。 \\
\opt{-target} & 指定仅提交 RedHawk 作业。 \\
\opt{-timeout} & 指定作业提交超时值。 \\
\opt{-submit\_command} & 指定 LSF、GRD、SGE 或定制作业提交程序的完整路径名以及其它所需选项，例如 \texttt{qsub} 或 \texttt{bsub}。 \\
\opt{-max\_cores} & 控制每个分布式 worker 作业可使用的核数。默认为 1。 \\
\bottomrule
\end{tabularx}

  \item 运行带下列选项的 \cmd{analyze\_rail}：

\begin{tabularx}{\textwidth}{@{}l X@{}}
\toprule
\textbf{选项} & \textbf{说明} \\
\midrule
\opt{-submit\_to\_other\_machines} & 启用向目标 farm 机器提交 RedHawk 作业。 \\
\opt{-multiple\_script\_files} \texttt{directory\_name script\_name} & 指定 RedHawk 脚本以及保存分析结果的目录。 \\
\bottomrule
\end{tabularx}
\end{enumerate}

\textbf{示例 51}
\begin{lstlisting}
%> rsh myMachine ls
test1
CSHRC
DESKTOP
> set_host_options "myMachine" -target RedHawk \
   -max_cores 2 -timeout 100
> analyze_rail -submit_to_other_machines \
   -voltage_drop static -nets {VDD VDDG VDDX VDDY VSS}
\end{lstlisting}
上例先检查机器 \texttt{myMachine} 是否支持 rsh 协议，再仅为 RedHawk Fusion 启用向该机器提交作业。超时值为 100，使用核数为 2。

\textbf{示例 52}
\begin{lstlisting}
> set_host_options -submit_protocol sge \
   -submit_command { qsub -P normal }
> analyze_rail -submit_to_other_machines -voltage_drop \
   static -nets {VDD VSS}
\end{lstlisting}
上例启用向 SGE farm 机器提交作业。

\textbf{示例 53}
\begin{lstlisting}
> set_host_options -submit_protocol lsf \
   -submit_command {bsub}
> analyze_rail -submit_to_other_machines \
   -multiple_script_files {{result1 RH1.tcl} {result2 RH2.tcl}}

> open_rail_result result1
> close_rail_result
> open_rail_result result2
\end{lstlisting}
上例用两个 RedHawk 脚本向 LSF farm 机器提交两个 RedHawk 作业，并将保存分析结果的目录指定为 \texttt{result1} 与 \texttt{result2}。

\textbf{示例 54}
\begin{lstlisting}
> set_app_options -name rail.license_checkout_mode \
   -value wait
> set_host_options -submit_protocol sge \
   -submit_command {qsub}
> analyze_rail -submit_to_other_machines \
   -multiple_script_files {{ result1 RH1.tcl} {result2 RH2.tcl} \
   {result3 RH3.tcl}}
\end{lstlisting}
上例向 SGE farm 机器提交三个 RedHawk 脚本，并在 \cmd{analyze\_rail} 运行期间等待可用许可证。

% ============================================================
\section{电压驱动的电源开关单元尺寸调整}
\subsection*{Voltage Driven Power Switch Cell Sizing}

Fusion Compiler 工具可调整电源开关单元尺寸，以降低总功耗并缓解设计中的压降。电源开关单元具有不同阈值电压（VT），分为 high、standard、low 与 ultra-low。高阈值电压的电源单元功耗与漏电流较低，但时序较慢、驱动强度较小；低阈值电压的电源单元时序较快，但功耗与漏电流较高。通过电源开关单元尺寸调整，工具支持不同驱动强度与阈值电压，以降低设计中的 IR（电压）降。

\begin{noteBox}
当工具切换电源单元（交换）时，只能交换 footprint 相同但阈值电压不同的单元。这意味着若单元尺寸与 pin 形状不同，则无法交换。
\end{noteBox}

要执行电源开关单元交换，可使用 \cmd{size\_power\_switches} 设置最大 IR drop 目标。交换电源单元以满足该目标时，工具尝试寻找 VT 更高但 $R_{\mathrm{ON}}$ 更低的另一单元，从而降低 IR drop。可用 \cmd{set\_power\_switch\_resistance} 设置工具在交换时考虑的 ON 电阻（$R_{\mathrm{ON}}$）值。基于 IR drop 目标与电源单元设定的 $R_{\mathrm{ON}}$，工具尝试替换电源单元以满足目标压降。

\begin{noteBox}
若工具找不到匹配目标压降与 $R_{\mathrm{ON}}$ 的合适电源单元，则跳过单元尺寸调整。
\end{noteBox}

要执行电压驱动的单元尺寸调整：
\begin{enumerate}
  \item 使用 RH Fusion 执行压降分析，获得由相关电源开关驱动的单元的最大压降。详见「执行压降分析」。
  \item 使用 \cmd{set\_power\_switch\_resistance} 设置单元切换时考虑的 $R_{\mathrm{ON}}$ 值限值。提供：
    \begin{itemize}
      \item Lib cell name — 库单元名。
      \item On resistance value — 单元尺寸调整时考虑的 $R_{\mathrm{ON}}$ 值限值。可从开关模型文件获取该信息。
    \end{itemize}
例如：
\begin{lstlisting}
fc_shell> set_power_switch_resistance \
 saed32hvt_pg_ss0p_v125c/HEAD2X2_HVT 0.015
\end{lstlisting}
  \item 运行 \cmd{size\_power\_switches} 执行单元尺寸切换。提供 \opt{-max\_irdrop} 的值，指定由开关驱动的标准单元在电源 net 上允许的最大有效压降。此为必填输入。例如：
\begin{lstlisting}
fc_shell> size_power_switches -max_irdop 0.05
\end{lstlisting}
关于进一步细化命令行为的其它非必填输入，见 \cmd{size\_power\_switches} 的 man page。
\end{enumerate}

工具还提供下列命令以进一步协助电源开关电阻：
\begin{itemize}
  \item \cmd{get\_power\_switch\_resistance} — 获取电源开关单元的 $R_{\mathrm{ON}}$ 值。
  \item \cmd{reset\_power\_switch\_resistance} — 重置为电源开关单元指定的 $R_{\mathrm{ON}}$ 值。
  \item \cmd{report\_power\_switch\_resistance} — 报告所有已设置 $R_{\mathrm{ON}}$ 的电源开关库单元。使用 \opt{-verbose} 可显示 footprint 相同的单元组、$R_{\mathrm{ON}}$ 值以及设计中例化的单元数。使用 \opt{-verbose} 时还可使用 \opt{-same\_vt}，报告当前库中 VT 相同但驱动强度不同的电源开关单元及其 $R_{\mathrm{ON}}$。
\end{itemize}
